% 全国大学生数学建模竞赛论文骨架（高教社杯/华数杯）
% 电子版：第一页为摘要页，不含承诺书与编号页。
% 正文约 20 页，附录代码不超过 25 页。
% 用 xelatex 编译：python scripts/compile_latex.py main.tex
\documentclass[withoutpreface,bwprint]{cumcmthesis}

\usepackage{amsmath,amssymb}
\usepackage{graphicx}
\usepackage{booktabs}
\usepackage{float}
\usepackage{algorithm}
\usepackage{algpseudocode}
\usepackage{listings}
\lstset{
  language=Python,
  basicstyle=\ttfamily\footnotesize,
  breaklines=true,
  frame=single,
  numbers=left,
  tabsize=4,
}

\title{题目名称}
\tihao{A}
\baominghao{参赛队号}
\schoolname{学校名称}
\membera{队员一}
\memberb{队员二}
\memberc{队员三}
\supervisor{指导教师}
\yearinput{2025}
\monthinput{09}
\dayinput{07}

\begin{document}

\maketitle

\begin{abstract}
摘要写满一页但不超页。先一句话写问题与方法，再按问给方法和数值结果，最后写创新点与检验结论。不要机械堆"针对问题一……问题二……"的排比开头。

\keywords{关键词1；关键词2；关键词3}
\end{abstract}

\tableofcontents
\newpage

\section{问题的背景与重述}

\subsection{问题背景}

\subsection{问题重述}

\section{问题分析}

按问分节，每问写三块：要求、难点、思路。每问配一张流程图。

\subsection{问题一的分析}

要求：……

难点：……

为此，本文……

\begin{figure}[H]
\centering
% 流程图占位，用 TikZ 或脚本生成
\caption{问题一流程分析图}
\end{figure}

\subsection{问题二的分析}

\subsection{问题三的分析}

\subsection{本文的主要工作与创新}

超出题目本身的内容放在这里或"模型的拓展"节，是争奖的关键。

\section{模型假设}

每条假设附一句合理性说明。

\begin{enumerate}
  \item 假设一。该假设的合理性：……
  \item 假设二。该假设的合理性：……
\end{enumerate}

\section{符号说明}

\begin{table}[H]
\centering
\begin{tabular}{cll}
\toprule
符号 & 含义 & 单位 \\
\midrule
$x$ & 示例 & \\
\bottomrule
\end{tabular}
\end{table}

\section{模型的建立与求解}

公式全部编号，推导写完整步骤。

\subsection{问题一}

\subsubsection{模型的建立}

\begin{equation}
E = mc^2 .
\end{equation}

\subsubsection{模型的求解}

\begin{algorithm}[H]
\caption{求解算法}
\begin{algorithmic}[1]
\State 第一步；
\State 第二步；
\end{algorithmic}
\end{algorithm}

\subsection{问题二}

\subsection{问题三}

\section{模型的拓展与深入分析}

方法理论分析、判据标定、不确定性传播、推广与设计准则。

\section{模型的检验}

合成数据验证表、一致性检验、敏感性分析。

\section{模型的评价}

优缺点与推广。

\section{参考文献}

\begin{thebibliography}{9}
\bibitem{example2026} 作者. 文献标题. 出版信息, 2026.
\end{thebibliography}

\appendix
\section{附录：完整代码}

代码文件清单与说明，结果文件说明。

\subsection{main.py}
\lstinputlisting{../codes/main.py}

\end{document}
